\section{Zusammenfassung}
In diesem Versuch wurden die Hysteresekurven von drei verschiedenen Ferromagnetika aufgenommen und
daraus die magnetischen Kenngrößen Remanenz $B_r$, Koerzitivfeldstärke $H_c$ und Sättigungsflussdichte $B_s$ bestimmt.
Die untersuchten Materialien waren zwei Nickel-Eisen-Legierungen (PERMENORM 5000 H2 und PERMENORM 5000 Z)
sowie ein Flachstahlkern. \\
PERMENORM 5000 H2 zeigte mit einer Remanenz von $B_r = 0.637 \pm 0.038 \, T$ und einer Koerzitivfeldstärke
von $H_c = 5.87 \pm 0.15 \, A/m$ die charakteristischen Eigenschaften eines weichmagnetischen Werkstoffes.
Die Sättigungsflussdichte wurde zu $B_s = 1.583 \pm 0.075 \, T$ bestimmt. 
Zudem wurde der Ummagnetisierungsverlust auf bestimmt: $\textcolor{red}{\omega}_{verlust, H2} = 3.88 \pm 0.68 \, \frac{mJ}{kg}$\\
Die Nickel-Eisen-Legierung PERMENORM 5000 Z zeigte mit $B_r = 1.389 \pm 0.012 \, T$ und
$H_c = 14.41 \pm 0.22 \, A/m$ ebenfalls weichmagnetische Eigenschaften, jedoch mit einer etwas
höheren Koerzitivfeldstärke als PERMENORM 5000 H2. Die Sättigungsflussdichte wurde zu
$B_s = 1.554 \pm 0.074 \, T$ bestimmt. \\
Der Flachstahlkern zeigte mit \textcolor{red}{$B_r = 0.5261 \pm 0.00018 \, T$ und $H_c = 344.9 \pm 1.9 \, A/m$}
deutlich andere magnetische Eigenschaften als die beiden Nickel-Eisen-Legierungen. Die
geringe Remanenz und die hohe Koerzitivfeldstärke deuten auf einen hartmagnetischen Werkstoff hin.
Die Sättigungsflussdichte konnte in diesem Fall nicht bestimmt werden, da die maximale
angelegte Feldstärke nicht ausreichte, um den Kern in die Sättigung zu bringen. 
Ebenso wurde auch hier der Ummagnetisierungsverlust bestimmt: $\textcolor{red}{\omega}_{verlust, Eisen} = \textcolor{red}{0.457 \pm 0.025} \, \frac{J}{kg}$\\
Zusätzlich wurde die Irreversibilität der Magnetisierung eines Eisenkerns durch Richtungsänderung
des Magnetfeldes untersucht, sowie die Entmagnetisierung durch sukzessive Verringerung
des maximalen Magnetfeldes nach jeder Umpolung demonstriert. \\
\textcolor{red}{Die angegebenen Unsicherheiten ergeben sich überwiegend aus dem
graphischen Ablesen der Hysteresekurven. Sowohl die Bestimmung der
Remanenz \( B_r \) als auch der Koerzitivfeldstärke \( H_c \)
erfolgte durch visuelles Identifizieren der Schnittpunkte mit den
Achsen. Aufgrund der endlichen Linienbreite der aufgezeichneten
Kurven, der begrenzten Auflösung der Anzeigeinstrumente sowie
möglicher Nullpunktverschiebungen ist diese Methode mit einer
systematischen Unsicherheit behaftet. 
Auch die Bestimmung der Sättigungsflussdichte \( B_s \) ist
unsicher, da der Übergang in die Sättigung keinen eindeutig
definierten Punkt besitzt, sondern sich nur asymptotisch annähert.
Die Berechnung der Ummagnetisierungsverluste über die eingeschlossene
Fläche der Hystereseschleife reagiert zudem empfindlich auf kleine
Verzerrungen oder Verschiebungen der Kurvenform.\\
Neben den Ableseunsicherheiten tragen insbesondere systematische
Effekte zur Gesamtabweichung der Ergebnisse bei. Während der
Messung kann es zu einem Drift des Integrators oder zu
Nullpunktverschiebungen in der Messverstärkerkette kommen,
wodurch sich die gesamte Hysteresekurve langsam verschiebt.
Eine solche Drift beeinflusst direkt die Bestimmung von
\( B_r \) und \( H_c \) und kann insbesondere bei kleinen
Remanenzwerten zu großen relativen Fehlern führen.
Ebenso wirken sich Unsicherheiten in den Spulenkonstanten,
in der Kalibrierung der Feldstärke sowie Schwankungen der
angelegten Stromstärke systematisch auf die berechneten
Feld- und Flussdichtewerte aus. Im Vergleich zu rein
statistischen Schwankungen dürften diese systematischen
Beiträge dominieren, da sie alle Messpunkte in gleicher
Weise verschieben und somit die gesamte Kurvenform beeinflussen.}
Des weiteren wurde das Verhalten von Holz unter Einfluss eines Magnetfeldes qualitativ untersucht.
Es zeigte sich, dass Holz diamagnetische Eigenschaften besitzt, da es in einem inhomogenen
Magnetfeld eine Kraft in Richtung der geringeren Feldstärke erfährt.
Die relative Permeabilität von Holz wurde zu $\mu_r = 4.571 \pm 0.082$ abgeschätzt. Dies ist nicht mit dem Literaturwert von
$\mu_{r,lit} < 1$ \cite{Holz} vereinbar, was auf die ungenaue Messmethode und auf die geringe Antwort von Holz
auf ein äußeres Magnetfeld zurückzuführen ist. 
\textcolor{red}{Die starke Abweichung vom Literaturwert verdeutlicht, dass die
verwendete Messmethode für die quantitative Bestimmung der
relativen Permeabilität diamagnetischer Materialien nur eingeschränkt
geeignet ist. Diamagnetismus ist ein sehr schwacher Effekt
(\( \mu_r \lesssim 1 \)), sodass bereits geringe systematische
Fehler — beispielsweise durch Inhomogenitäten des Magnetfeldes,
Unsicherheiten in der Kraft- oder Wegmessung sowie Drift der
Anzeige — zu einer erheblichen relativen Abweichung führen.
Die beobachtete Diskrepanz ist daher primär auf systematische
Messunsicherheiten und nicht auf statistische Streuung
zurückzuführen.}
